\documentclass[12pt,a4paper]{article}

%------------------------------------------------
% PACKAGES
%------------------------------------------------
\usepackage[utf8]{inputenc}
\usepackage[T1]{fontenc}
\usepackage{lmodern}
\usepackage{longtable}
\usepackage{amsmath, amssymb}
\usepackage{graphicx}
\usepackage{booktabs}
\usepackage{geometry}
\usepackage{natbib}
\usepackage{setspace}
\usepackage{multirow}
\usepackage{multirow}
\usepackage{lscape}
\usepackage{tabularx}
\usepackage{booktabs}
\usepackage{multirow}
\usepackage{array}
\usepackage{hyperref}
\geometry{margin=2.5cm}
\onehalfspacing

%------------------------------------------------
% TITLE
%------------------------------------------------
\title{What Makes a University Sustainable in the Eyes of its Core Community?\\
\large Sustainability Expectation--Perception Gap}

\author{Author Name(s)\\
Institution\\
Email}

\date{}

\begin{document}

\maketitle

%------------------------------------------------
% ABSTRACT
%------------------------------------------------


\textbf{Purpose} -- This study adopts a behavioral, stakeholder-centered approach to identify the key dimensions of university sustainability valued by internal stakeholders (students, faculty members and administrative staff). It examines how they interpret and prioritize environmental, social, and economic dimensions, assesses gaps between expectations and perceived initiatives, and proposes a composite score to support the University Sustainability Performance Index (USPI).

\textbf{Design/methodology/approach} -- This qualitative case study at a single university used semi-structured interviews with students, faculty, and staff to explore perceptions of environmental, social, and economic sustainability. Thematic analysis identified key dimensions and gaps between initiatives and expectations, forming the basis for a composite score underpinning the University Sustainability Performance Index (USPI).

\textbf{Findings} -- The results reveal gaps between stakeholder perceptions and expectations. Socially, student well-being and community engagement need stronger structuring. Environmental initiatives are recognized, yet sustainable mobility and awareness require better formalization. Economic gaps remain significant, especially in energy optimization, financial governance, and sustainable investment.

\textbf{Research limitations} -- The study focuses on internal stakeholders in a specific context. Future research could include external stakeholders and cross-institutional comparisons. It extends stakeholder theory by applying a behavioral lens and frames sustainability as a co-created outcome.

\textbf{Practical implications} -- The USPI provides university leaders with a transparent, stakeholder-driven decision-support tool to assess sustainability performance, prioritize strategic actions and allocate resources effectively. It facilitates evidence-based governance aligned with stakeholder expectations.

\textbf{Social implications} -- By aligning institutional sustainability practices with stakeholder-defined priorities, universities can enhance their societal resonance, strengthen trust and contribute more effectively to the Sustainable Development Goals (SDGs), particularly SDGs~4, 8, 9 and 17.

\textbf{Originality/value} -- This study moves beyond traditional operational sustainability metrics by adopting a behavioural and stakeholder-centred approach. It introduces a composite index grounded in stakeholder preferences, offering a novel answer to the question of what defines a sustainable university from the perspective of its core academic community.


\noindent\textbf{Keywords} -- University sustainability; Sustainability Expectation--Perception Gap; SDGs; USI; University sustainability index.

\noindent\textbf{Article Type} -- Research paper

\section{Introduction}

Sustainability in higher education institutions (HEIs) has become a strategic and societal imperative in response to growing environmental, social, economic and governance challenges worldwide. Universities are no longer perceived solely as knowledge-producing entities; they are increasingly expected to act as transformative agents capable of shaping sustainable development trajectories. In alignment with the United Nations Sustainable Development Goals (SDGs), sustainability in HEIs extends beyond environmental stewardship to encompass social equity, economic viability, institutional governance and academic excellence.

Despite the proliferation of sustainability rankings, reporting frameworks and performance indicators, existing assessment tools often remain predominantly operational and managerial in nature. They typically rely on externally defined metrics and institutional reporting mechanisms, while giving limited attention to how sustainability is perceived and prioritized by the university’s own academic community. This represents a significant conceptual gap. Sustainability is not merely a technical or administrative construct; it is also a socially constructed and behaviorally interpreted phenomenon shaped by stakeholder perceptions, expectations and value systems.

Stakeholder theory \cite{Freeman1984} posits that organizational sustainability depends on the integration of diverse stakeholder interests. Within HEIs, internal stakeholders, particularly students and faculty, play a central role in shaping the institution’s sustainability identity. Students seek meaningful education, employability and ethical alignment with societal challenges, while faculty members emphasize academic integrity, research impact and institutional mission. Yet, little empirical research has systematically examined how these core actors interpret sustainability, how they prioritize its dimensions, and how they evaluate institutional performance relative to their expectations.

Addressing this gap is essential for developing assessment tools that genuinely reflect institutional realities rather than externally imposed benchmarks. If sustainability indices are to guide strategic decision-making effectively, they must be grounded in stakeholder-defined priorities. In this context, understanding what makes a university ``sustainable'' in the eyes of its core community becomes a critical research question.

This study adopts a behavioural and stakeholder-centred perspective to analyse the sustainability preferences of students, faculty members, and administrative staff within a case-study university. Through qualitative inquiry, it explores how these actors interpret environmental, social and economic sustainability dimensions, how they hierarchize them, and how they perceive institutional performance. By identifying expectation--performance gaps, the study provides an empirical foundation for proposing a composite sustainability score aligned with stakeholder-defined priorities.

By doing so, this research contributes to the broader University Sustainability Performance Index (USPI) project by grounding the index in stakeholder perceptions rather than solely in technical criteria. The study advances sustainability research in higher education by integrating stakeholder theory with a behavioural lens and by conceptualizing sustainability as a co-created outcome within the university ecosystem. Ultimately, it seeks to provide both theoretical insight and practical guidance for developing more legitimate, inclusive and socially resonant sustainability assessment frameworks.

\section{Literature Review}

\subsection{Stakeholder Theory and University Sustainability}

Stakeholder theory provides a foundational lens for understanding sustainability in higher education institutions (HEIs). Originally developed in strategic management \citep{Freeman1984}, the theory argues that organizational legitimacy and long-term performance depend on the ability to recognize, balance, and integrate the interests of multiple stakeholder groups rather than privileging a single dominant constituency. Organizations are embedded in networks of relationships, and value is generated through ongoing interactions among actors with diverse expectations.

\citet{DonaldsonPreston1995} further conceptualized stakeholder theory along descriptive, instrumental, and normative dimensions. From a normative perspective, stakeholders possess intrinsic legitimacy; their interests are not merely strategic variables but morally relevant claims that organizations are obliged to consider. This normative foundation is particularly relevant in the context of universities, which are mission-driven institutions with societal responsibilities extending beyond economic performance.

\citet{Mitchell1997} introduced the stakeholder salience framework, emphasizing that stakeholders differ in terms of power, legitimacy, and urgency. Within HEIs, a wide range of stakeholders coexist, including students, faculty members, administrative staff, employers, accreditation bodies, local communities, and policymakers. However, internal stakeholders, particularly students and faculty, occupy a structurally central position. They are not only recipients of sustainability policies but active agents who interpret, enact, and sometimes contest institutional sustainability agendas.

More recent developments in stakeholder theory emphasize relationality and value co-creation \citep{Freeman2010}. From this perspective, sustainability is not merely implemented through top-down governance mechanisms or infrastructure investments; it emerges from shared meaning-making processes and collective engagement. Universities can thus be understood as ecosystems in which sustainability outcomes are co-produced through interactions between institutional strategies and stakeholder perceptions.

While stakeholder theory provides a strong normative and relational foundation for understanding sustainability in higher education, the empirical literature on campus sustainability has developed along several thematic trajectories.

Higher Education Institutions (HEIs) are increasingly recognized as pivotal actors in fostering sustainable development and addressing global environmental challenges. Often described as ``small cities'' due to their significant populations and complex operations \citep{Alshuwaikhat2008, Velazquez2006, MinAllah2020, Murshed2020}, universities wield substantial direct and indirect impacts on the environment through their energy consumption, material use, waste generation, and daily activities. Consequently, the concept of a ``green campus'' has emerged as a comprehensive approach to integrating sustainability principles across physical infrastructure, academic curricula, research, and institutional governance \citep{Disterheft2013, Sugiarto2022}.

This literature review synthesizes recent research on green campus initiatives, grouping papers by key themes. It examines various approaches to conceptualizing and assessing campus sustainability, details diverse operational strategies, explores educational and behavioral aspects, and highlights the crucial role of governance and policy. The review also identifies prominent methodological trends and concludes with a summary of key insights and directions for future research.

\subsection{Conceptualization, Frameworks, and Assessment Tools for Green Campuses}

A significant portion of the literature focuses on defining, conceptualizing, and assessing sustainability within HEIs. Early work proposed comprehensive managerial models for a ``sustainable university'' \citep{Velazquez2006} and integrated approaches to campus sustainability \citep{Alshuwaikhat2008}. The evolving nature of ``campus sustainability'' is often understood in a broad sense, encompassing physical, educational, and institutional dimensions \citep{Disterheft2013}. Researchers note a shift from simplistic ``going green'' trends to a more holistic, systemic understanding that integrates environmental, social, and economic dimensions \citep{Alshuwaikhat2016, Sugiarto2022, Popowska2024}.

The development and application of assessment tools and ranking systems are central to measuring progress and benchmarking sustainability efforts. Prominent examples include the UI GreenMetric World University Rankings \citep{Suwartha2013, MunozSuarez2020, Alawneh2021, Safarkhani2022, Fatah2025} and the Sustainability Tracking, Assessment \& Rating System (STARS) \citep{Urbanski2015, Kirrane2020, Zhu2021, WashingtonOttombre2024}. Other frameworks like AISHE (Auditing Instrument for Sustainability in Higher Education) \citep{Lambrechts2013, Lambrechts2015}, SAQ (Sustainability Assessment Questionnaire) \citep{Alshuwaikhat2016}, ISCN (International Sustainable Campus Network) \citep{Sonetti2016}, CANSEC (Carbon Neutrality and Sustainability in Educational Campuses) \citep{Jain2017}, and PICSOU (Performance Indicators for Core Sustainability Objectives of Universities) \citep{Jiang2023} are also utilized.

Challenges with these tools are frequently discussed, including their potential bias towards energy efficiency indicators \citep{Sonetti2016}, lack of standardization, and difficulties in adapting to context-specific issues, especially in developing countries \citep{Velasco2018, Zhu2021, Dawodu2024}. There is also a recognized need for more robust tools that encompass social and economic dimensions holistically \citep{Sonetti2016, Zhang2025} and ensure meaningful international comparisons \citep{Horan2021}. The integration of such tools into strategic planning remains an ongoing area of research \citep{Chen2019, Horan2021, Alalawi2024}.

\subsection{Operational Sustainability Initiatives}

A significant body of research focuses on concrete actions and initiatives to reduce the environmental footprint of campus operations.

\subsubsection{Energy Efficiency \& Renewable Energy}

This is a dominant theme, with studies analyzing energy consumption patterns \citep{Maistry2016}, proposing energy efficiency plans for buildings \citep{Soares2015, Pini2019, Ahmed2025, Laporte2025, Movlyanov2025}, and optimizing energy systems \citep{Amaral2020, Zhu2021}. Many papers investigate the feasibility and implementation of renewable energy sources, particularly solar PV systems, for on-campus electricity generation \citep{Muller2023, Diniz2024, Ayadi2025, Nayan2023}. These initiatives often aim for carbon emission reduction \citep{Odabasi2023, Mohamad2024}, utilizing technologies like motion sensors \citep{Riyanto2018} and smart grids \citep{Sharma2016}. Some studies explicitly link energy saving to a broader sustainability impact \citep{Zhu2021}.

\subsubsection{Waste Management \& Circularity}

Research highlights the importance of effective waste management, including 3R (Reduce, Reuse, Recycle) initiatives. Studies assess food waste composting \citep{AlAomar2022, Rupani2023}, plastic waste reduction through bottle refill stations \citep{Pujiati2024}, and the institutionalization of waste minimization \citep{Zen2016, Agustina2023}. The concept of circularity is explored, particularly in managing IT material flows \citep{Leclerc2022}. Waste characterization studies are key to developing tailored management strategies \citep{Hindiyeh2022, Devec2025}.

\subsubsection{Water Management}

Initiatives focus on water conservation, rainwater harvesting systems and smart water management using IoT \citep{Nadzri2024, Siddamal2022}. Water quality assessment and wastewater reuse are also explored \citep{Lopes2017, DeWata2025}.

\subsubsection{Transportation \& Mobility}

Reducing the carbon footprint from commuting is a significant concern. Studies analyze travel patterns, develop sustainable mobility action plans \citep{Ibrahim2022}, and promote cycling and public transport \citep{MateoBabiano2020, Ibrahim2022, Saini2023}. The effectiveness of bus operations is also optimized \citep{Johar2025}.

\subsubsection{Green Buildings \& Infrastructure}

Research investigates green building principles, energy-efficient refurbishment and passive design strategies for thermal comfort. The role of green infrastructure in sustainability goals is also examined \citep{Tao2019, Pini2019, Kanchana2025}.

\subsubsection{Land Use \& Biodiversity}

Studies address optimizing green spaces for mental health and well-being \citep{Mustari2025}, carbon sequestration by urban trees \citep{Yasin2024, Lee2025}, and developing campus botanical gardens \citep{Hapsari2025, Mustari2025}.

\subsection{Education, Awareness \& Engagement}

HEIs have a crucial role in shaping a sustainability-conscious society through education and fostering pro-environmental behaviors.

\subsubsection{Curriculum Integration \& Education for Sustainable Development (ESD)}

Many papers discuss integrating sustainability into curricula \citep{Mendez2025, Zhu2025, Builes2024, Orhan2021}. This includes developing specific courses or modules \citep{Nagatomo2023}, using living laboratories for experiential learning \citep{Gomez2021}, and promoting interdisciplinary approaches \citep{Nagatomo2023}. Challenges include addressing misinformation about climate change \citep{Hess2018} and ensuring holistic integration beyond technical knowledge \citep{Gamage2022}. The role of ESD in different contexts, including developing and emerging countries, is also highlighted.

\subsubsection{Student/Staff Awareness \& Behavior Change}

Fostering awareness and encouraging pro-environmental behaviors among students and staff is critical \citep{Setiawati2021, Alalawi2024}. Studies explore factors influencing recycling behavior and energy conservation \citep{Thoo2022, Ishak2016}. The impact of social media in disseminating awareness is also noted \citep{Hamid2017}. Research often investigates the gap between awareness/expectations and actual behavior or university performance \citep{Dagiliute2018, Zhao2023, Yasin2024}. Studies also examine interventions to promote sustainable actions, such as ``green office initiatives'' \citep{Zen2016} or ``energy challenges'' \citep{Boulton2017}. The ``heartware'' approach, focusing on internal motivation and collective efforts, is highlighted as a driver for long-term sustainability \citep{Mohamad2018}.

\subsubsection{Community Engagement \& Social Responsibility}

Universities extend their sustainability efforts to the broader community through outreach programs and by acting as role models \citep{Fatah2025}. University Social Responsibility (USR) is a concept gaining traction, emphasizing universities' impacts on society and the environment beyond their immediate operations \citep{Gomez2018, AriasValle2024}. Studies explore student leadership \citep{Mohamed2025}, faculty involvement \citep{Alali2024}, and the potential for campus community gardens to foster direct engagement with sustainability \citep{DeYoung2016, Dantas2018}.

\subsection{Governance, Policy, and Organizational Change}

Effective governance and strategic policy are crucial for embedding sustainability within HEIs.

\subsubsection{Strategic Planning \& Institutionalization}

Papers emphasize the need for clear institutional frameworks \citep{Alshuwaikhat2008, Abubakar2020}, strategic sustainability plans \citep{White2014, Romero2017, Eby2023}, and masterplans that define policies, goals, and actions \citep{Romero2017}. The integration of sustainability into institutional policies is seen as essential for long-term commitment \citep{DaSilva2023, Alalawi2024}. Shared governance and inclusive decision-making processes are identified as critical factors for success \citep{DjukicMin2025}.

\subsubsection{Barriers and Drivers to Implementation}

Research identifies various barriers, including lack of awareness, knowledge, financial constraints, resistance to change, inefficient communication, and lack of external incentives \citep{Alshuwaiwat2008, Alalawi2024, Herzanita2024, Geh2024}. Conversely, drivers include governmental pressure, financial savings, institutional leadership, student interest, and collaborative efforts \citep{Geh2023}. The ``fear of not flying'' due to career progression pressures can also hinder sustainable academic travel \citep{NurseyBray2019}.

\subsubsection{Organizational Learning \& Reporting}

Sustainability reporting is recognized as a tool for increasing organizational learning and strengthening sustainability efforts \citep{WashingtonOttombre2024}. The use of metrics and reporting frameworks helps track progress, but deep organizational learning is often elusive \citep{WashingtonOttombre2024, WashingtonOttombreBigalke2018, Rosa2024}.

\subsection{Smart Campus and Digitalization for Sustainability}

An emerging theme is the integration of smart technologies and digitalization to enhance campus sustainability.

\subsubsection{Smart Campus Concepts and Applications}

The ``smart campus'' concept leverages smart technologies with physical infrastructure for improved services, decision-making, and sustainability \citep{MinAllah2020, Ahmed2023, Alalawi2024, Sugiarto2022}. This includes smart microgrids \citep{MinAllah2020}, intelligent monitoring and visualization systems \citep{Ceccarini2021}, and IoT platforms for managing energy, water, and security \citep{Subbarao2019, Sanat2024, Alalawi2024}.

\subsubsection{Digital Tools for Engagement and Efficiency}

Social media is identified as a potent tool for disseminating environmental sustainability awareness among students and staff \citep{Hamid2017}. Gamified Human-Machine Interfaces (HMI) can engage users in energy reduction activities \citep{Mendez2021}. The research on ``green campus'' underscores the multifaceted and evolving role of HEIs in advancing sustainability. Universities are increasingly seen as living laboratories and models for sustainable development, influencing society through their operations, teaching, research, and outreach. Key themes reveal a progression from focusing primarily on environmental management (e.g., energy, waste, water) to embracing holistic sustainability that integrates social equity, economic viability, and institutional governance. Significant progress has been made in developing assessment tools and ranking systems to measure sustainability performance, although challenges remain in standardization, context-specificity, and comprehensive coverage of all sustainability dimensions. Operational initiatives, particularly in energy efficiency and renewable energy adoption, waste management, and sustainable transportation, are widespread. However, achieving deep, transformative change requires addressing behavioural aspects, fostering environmental awareness among all stakeholders, and embedding sustainability into the core curriculum.

\begin{longtable}{p{3cm} p{4cm} p{3cm} p{3cm} p{3cm}}
\caption{Summary of Analysed Studies} \\

\hline
\textbf{Theme} & \textbf{Key references (non-exhaustive)} & 
\textbf{Dominant methodological approaches} & 
\textbf{Convergent findings} & 
\textbf{Persistent limitations / research gaps} \\
\hline
\endfirsthead

\hline
\textbf{Theme} & \textbf{Key references (non-exhaustive)} & 
\textbf{Dominant methodological approaches} & 
\textbf{Convergent findings} & 
\textbf{Persistent limitations / research gaps} \\
\hline
\endhead

Conceptual frameworks \& sustainable campus models &
Velazquez et al. (2006); Alshuwaikhat \& Abubakar (2008); 
Disterheft et al. (2013, 2015); Sonetti et al. (2016); 
Sugiarto et al. (2022); Popowska \& Sady (2024) &
Conceptual papers, systemic models, integrative frameworks &
Holistic vision integrating operations, teaching, research, and governance &
Limited empirical validation and operationalization challenges \\

Assessment tools, indicators \& rankings &
Suwartha \& Sari (2013 UI GreenMetric); Urbanski \& Leal (2015); 
Lambrechts \& Ceulemans (2013 – AISHE); Sonetti et al. (2016); 
Zhu \& Dewancker (2021); Horan \& O’Regan (2021); 
Jiang \& Kurnitski (2023) &
Benchmarking systems, audits, questionnaires, scoring indices &
Enable reporting, comparison across universities, and strategic planning &
Energy/carbon bias, limited contextual adaptability, weak standardization in developing countries \\

Operational sustainability initiatives (energy, water, waste, mobility) &
Soares et al. (2015); Amaral et al. (2020); Zen et al. (2016); 
Ahmed et al. (2025) &
Energy simulations, LCA, technical audits, case studies &
Measurable reductions in consumption and emissions; strong technical focus &
Fragmented initiatives, weak integration of social and behavioral dimensions \\

Education for Sustainable Development (ESD) \& curriculum integration &
Burandt \& Barth (2010); Disterheft et al. (2013); 
Mendez et al. (2025) &
Pedagogical case studies, living labs, interdisciplinary approaches &
Development of sustainability competencies and experiential learning &
Limited assessment of long-term behavioral impacts \\

Awareness, behavior \& stakeholder engagement &
Dagiliūtė et al. (2018); Setiawati \& Maulana (2021); 
Zhao \& Zou (2023); Hamid et al. (2017) &
Surveys, statistical analyses, behavioral interventions &
High environmental awareness but persistent attitude–behavior gap &
Insufficient understanding of underlying motivational mechanisms \\

Governance, strategy \& organizational change &
White (2014); Romero \& Mulfarth (2017); 
Abubakar et al. (2020); Rosa et al. (2024); 
Washington-Ottombre (2024) &
Policy analysis, reporting studies, institutional case studies &
Leadership and strategic planning are critical drivers of sustainability performance &
Slow cultural change and limited internal coordination \\

University social responsibility \& community engagement &
Gómez et al. (2021); Arias-Valle \& Marimón (2024) &
Qualitative evaluations, impact studies &
Universities increasingly act as societal sustainability leaders beyond campus boundaries &
Lack of standardized metrics for social impact assessment \\

Smart campus \& digitalization for sustainability &
Min-Allah \& Alrashed (2020); Subbarao et al. (2019); 
Ceccarini et al. (2021); Ahmed et al. (2023) &
IoT systems, smart grids, AI-based monitoring platforms &
Real-time resource optimization and improved decision-making &
Limited longitudinal evidence and adoption/acceptance challenges \\

\hline
\end{longtable}

\noindent\textit{Source: Authors}


Despite this theoretical relevance, empirical sustainability measurement in higher education remains largely detached from stakeholder-centered approaches. Many indices prioritize quantifiable operational metrics, energy efficiency, waste management, infrastructure performance, while underrepresenting how sustainability is cognitively framed and hierarchized by internal actors. This disconnect creates a potential legitimacy gap between institutional sustainability strategies and stakeholder-defined priorities.

By mobilizing stakeholder theory within a behavioral framework, the present study positions sustainability as a socially negotiated construct shaped by expectation, perception dynamics. Understanding how core stakeholders define, prioritize and evaluate sustainability dimensions becomes essential for designing assessment tools that are not only technically robust but institutionally legitimate.

\section{Methodology}

\subsection{Research Design}

This study adopts an exploratory qualitative case study design aimed at capturing stakeholders’ perceptions of university sustainability within a specific higher education institution. A single-case approach was selected to allow an in-depth understanding of how sustainability is interpreted, prioritized and evaluated by core internal stakeholders. This design is particularly appropriate when investigating complex, context-dependent phenomena such as institutional sustainability, which involves multiple dimensions and heterogeneous actor perspectives.

The research is grounded in stakeholder theory \citep{Freeman1984} and adopts a behavioural perspective to examine how individual actors construct meaning around environmental, social and economic sustainability objectives within the university context.

\subsection{Data Collection}

Data were collected through semi-structured interviews conducted with key internal stakeholders, including: 26 Students, 18 Faculty members and 10 Administrative staff.

Participants were selected using purposive sampling to ensure the inclusion of stakeholders directly involved in sustainability-related activities and institutional processes. Specifically, students engaged in sustainability projects, faculty members involved in governance or research initiatives, and administrative staff responsible for institutional policies were intentionally recruited. The objective was not statistical representativeness but rather informational richness, enabling an in-depth exploration of diverse yet relevant perspectives on university sustainability.

Interviews allowed participants to express their own understanding of sustainability and to articulate perceived gaps between institutional commitments and actual practices. The interview guide was structured around three main themes: Perceived essential dimensions of university sustainability, interpretation and prioritization of sustainability objectives, expectations and perceived institutional performance.

\subsection{Data Analysis}

The data were analysed using thematic analysis. The analysis followed a systematic multi-step process:

\begin{itemize}
    \item Transcription of interviews
    \item Open coding to identify first-order concepts emerging from stakeholder discourse
    \item Axial coding to group concepts into broader sustainability dimensions and SDGs
    \item Identification of recurring patterns related to expectations, prioritization and perceived performance gaps
\end{itemize}

This iterative process enabled the identification of stakeholder-defined sustainability dimensions and the evaluation criteria implicitly used by participants.

The analysis aimed not only to categorize sustainability dimensions (environmental, social, economic), but also to capture how stakeholders hierarchize these dimensions and attach value to them.





\subsection{Development of the Composite Score}

Based on qualitative analysis conducted with the university’s core stakeholders, sustainability emerges as a multidimensional latent construct encompassing economic, social, and environmental dimensions \cite{1,2}. Within higher education institutions, sustainability has increasingly been conceptualized as an integrative framework combining environmental responsibility, social inclusion, and economic viability \cite{3}. The findings indicate that sustainability cannot be reduced to the mere presence of technical initiatives or institutional policies; rather, it must incorporate how these initiatives are perceived and evaluated by stakeholders. In line with stakeholder theory \cite{4}, institutional sustainability must be understood through stakeholder expectations, perceptions, and legitimacy dynamics \cite{5}. In this sense, sustainability constitutes not only an operational outcome but also a perceptual and relational construct shaped by stakeholder interpretation.

The qualitative evidence revealed a systematic divergence between stakeholder perceptions of existing sustainability initiatives and their expectations regarding institutional sustainability performance. Although operational actions were identified across economic, social, and environmental domains, stakeholders frequently highlighted shortcomings related to visibility, institutionalization, and strategic coherence. This divergence reflects a structural misalignment between institutional sustainability efforts and stakeholder priorities—an aspect that traditional sustainability indicators often fail to capture adequately \cite{6,7}.

To address this limitation, the present study introduces the Expectation--Perception Gap as a central conceptual construct for sustainability assessment. The conceptual foundation of this gap-based logic draws on the service quality literature, particularly the SERVQUAL model \cite{8}, which conceptualizes performance evaluation as the difference between expectations and perceptions. The gap represents the difference between stakeholder expectations and their perceptions of implemented sustainability initiatives, thereby capturing the degree of alignment between institutional performance and stakeholder priorities. A smaller gap indicates stronger alignment and greater perceived sustainability effectiveness, whereas a larger gap signals structural incoherence.

Within this framework, sustainability is modeled as a hierarchical latent construct structured around three core dimensions: economic, social, and environmental sustainability. The modeling approach is grounded in structural equation modeling theory, which conceptualizes sustainability as a higher-order latent variable inferred from observed indicators \cite{9,10}. Each dimension is jointly shaped by stakeholder perception and stakeholder expectation, and their interaction determines the level of institutional alignment.

Building upon this gap-centered perspective, the University Sustainability Index (USI) is introduced as an integrative higher-order latent construct representing overall institutional sustainability. Rather than relying exclusively on operational indicators, the USI reflects the degree of alignment between institutional sustainability initiatives and stakeholder expectations. This formulation advances stakeholder-centered sustainability assessment in higher education and responds to recent calls for more integrative and perception-based sustainability measurement frameworks \cite{7}.

\section{Results}

This section presents the conceptualization of the University Expectation--Perception Gap, based on qualitative analyses with the university’s internal stakeholders. The findings reveal areas of partial alignment as well as notable divergences between stakeholder perceptions and institutional initiatives.

\subsection{Qualitative conceptualization of the University Expectation--Perception Gap}

The dimension-based analysis identifies three main types of gaps (economic, social, and environmental) reflecting varying levels of misalignment between the university’s perceived performance and stakeholders’ expectations.

\subsection{Economic Expectation--Perception Gap}

The analysis of perceptions in the economic dimension reveals a predominantly operational understanding of sustainability, structured around three main themes: (1) digitalization, (2) energy optimization, and (3) human capital development.

Digitalization emerges as the dominant reference point, primarily associated with paper reduction, the dematerialization of administrative and pedagogical procedures, and the implementation of technological tools such as smart classrooms. This reflects a resource-efficiency logic centered on tangible and immediately observable material savings, particularly emphasized by students. In contrast, development of human capital constitutes a more strategic dimension, less spontaneously identified and mainly highlighted by faculty members.

Stakeholders’ expectations emphasize resource optimization and human capital development, while also introducing two previously unrecognized themes: (4) Responsible Financial Governance and (5) Responsible and Sustainable Investment.

\subsubsection{Strategic Alignment: Digitalization}

Digitalization emerges as a central theme in perceived economic sustainability, mentioned by the majority of respondents. It mainly includes paper reduction, the digitalization of administrative procedures, document dematerialization, and connected classrooms, with paper reduction being the most visible and most frequently cited measure. Economic sustainability is therefore primarily understood in operational terms, through a logic of direct material resource savings. Students particularly emphasize this dimension, while faculty members and administrative staff adopt a more balanced perspective.

From a gap perspective, digitalization appears only in perceptions and not in expectations, indicating alignment between perceptions and expectations. Stakeholders acknowledge the existing institutional efforts in this area without expressing additional demands, suggesting that this dimension reflects perceived operational performance rather than a strategic shortfall.

\subsubsection{Alignment Reinforcement Gap: Energy Optimization}

Energy optimization is a central sustainability theme within the organization, aiming to reduce consumption while maintaining comfort and service levels. It is divided into two main sub-themes: resource optimization and energy control. Resource optimization, frequently highlighted, includes smart space management, efficient use of equipment and consumables, and waste reduction, reflecting attention to daily practices. Energy control involves regulating lighting, heating, and air conditioning, as well as using sensors and automatic systems to monitor and manage energy consumption.

Analysis of the interviews reveals a significant gap between perceived initiatives and expectations. Current actions are largely visible, whereas expectations are more technical, systematic, and ambitious. These include wider deployment of solar panels, greater use of renewable energy, building insulation, widespread use of LED lighting, smart energy management systems, regular energy audits, and measurable performance monitoring regarding energy control.

\begin{quote}
I would like the school to make greater investments in energy efficiency solutions, such as the installation of solar panels, LED lighting systems, and other technologies that reduce energy consumption while maintaining comfort and performance of the facilities (Respondent 38, Faculty members).
\end{quote}

\subsubsection{Support Gap: Human Capital Development}

Human capital development is the least perceived and spontaneously mentioned theme by respondents. Unlike more tangible themes, such as resource optimization, this dimension is less directly visible. Observations indicate that faculty members, particularly those involved in training programs, academic partnerships, and initiatives supporting professional integration, are the main reporters of existing efforts.

In this domain, perceptions and expectations are largely aligned. Stakeholders recognize initiatives such as the integration of CSR into courses, sustainable entrepreneurship projects, incorporation of SDGs into programs, the internal incubator, and training aimed at future-oriented careers. Their expectations focus on expanding responsible internships, strengthening incubators, and enhancing training in responsible management.

Overall, the gap in this area is low, indicating strong coherence between existing initiatives and stakeholder expectations. This alignment suggests that current efforts in human capital development effectively meet anticipated needs, demonstrating that the institution’s strategic and operational actions in this domain are largely on target.

\subsubsection{Institutional Gap: Responsible Financial Governance}

This theme primarily emerged at the level of stakeholders’ expectations, with 30 mentions, highlighting its strategic importance. Expectations in this domain are far more explicit than the perceived initiatives of the organization. Regarding budget transparency, this sub-theme is entirely absent from perceptions: no mechanisms for monitoring or reporting on budget usage were reported. Stakeholders express clear and ambitious expectations, including the establishment of a dedicated sustainability budget, institutionalized green funds, public financial reports, and quantifiable indicators to track the impact of sustainability actions.

\begin{quote}
I would like to see greater transparency in the use of budgets related to sustainability, so that everyone can clearly understand how financial resources are allocated and utilized, as well as the creation of dedicated funds for green projects (Respondent 18, Administrative staff).
\end{quote}

Similarly, for the communication of results, perceived actions remain largely invisible to stakeholders. They expect a well-defined communication strategy, the publication of an annual CSR report, and regular dissemination of sustainability outcomes.

The gap identified is therefore institutional in nature, reflecting a lack of formalization and structured processes for existing initiatives. Stakeholders’ expectations emphasize the need to establish robust institutional mechanisms that ensure financial transparency and effective communication of results, thereby strengthening the organization’s credibility and commitment to sustainability.

\subsubsection{Strategic Gap: Responsible Investment}

The theme of responsible and sustainable investment also emerged prominently in stakeholders’ expectations, with 24 mentions, underscoring its strategic significance. While some sustainability-related economic initiatives exist, such as resource optimization, digitization of processes, and a few student projects, perceived actions remain limited compared to stakeholder expectations. Two main sub-themes are identified: green budgets and ecological investment. Stakeholders express the need for institutionalized budgets dedicated to sustainability, enabling the planned and sustainable funding of green and responsible projects.

\begin{quote}
The school has already undertaken certain sustainability initiatives, but I would like to see greater transparency regarding the budgets allocated to sustainable development, in order to better understand how resources are used and to strengthen the credibility and impact of the actions undertaken. (Respondent 28, Student).
\end{quote}

Regarding ecological investment, perceived initiatives largely involve occasional energy savings and digitalization efforts. However, stakeholders expect a broader and more structured approach, including diversified investments in ecological solutions, responsible procurement practices, energy-efficient equipment, and regular monitoring of environmental impact.

The identified gap is therefore strategic, reflecting the need to strengthen the planning and formalization of sustainable investments. Current initiatives remain visible, whereas stakeholder expectations align with a comprehensive strategic vision encompassing resource allocation, systematic monitoring, planning, and continuous improvement of sustainable economic practices.

\renewcommand{\arraystretch}{1.2}
\footnotesize

\begin{longtable}{%
>{\raggedright\arraybackslash}p{2.4cm}
>{\raggedright\arraybackslash}p{2.7cm}
>{\raggedright\arraybackslash}p{3.4cm}
>{\raggedright\arraybackslash}p{3.6cm}
>{\raggedright\arraybackslash}p{2.6cm}
>{\centering\arraybackslash}p{1.2cm}
}

\caption{Economic Expectation--Perception Gap Analysis}
\label{tab:economic_gap} \\

\toprule
\textbf{Dimension} & 
\textbf{Sub-theme} & 
\textbf{Perception (existing initiatives)} & 
\textbf{Priority expectations} & 
\textbf{Nature of the gap} & 
\textbf{Level} \\
\midrule
\endfirsthead

\toprule
\textbf{Dimension} & 
\textbf{Sub-theme} & 
\textbf{Perception (existing initiatives)} & 
\textbf{Priority expectations} & 
\textbf{Nature of the gap} & 
\textbf{Level} \\
\midrule
\endhead

\midrule
\multicolumn{6}{r}{\textit{Continued on next page}} \\
\midrule
\endfoot

\bottomrule
\endlastfoot

Digitalization &
Paper reduction and digital processes &
Dematerialization of administrative procedures; connected classrooms; reduction of paper use &
No additional expectations explicitly expressed &
Strategic alignment &
None \\

\addlinespace

\multirow{2}{*}{Energy optimization}
& Resource optimization
& Limited solar panels; occasional efficiency initiatives
& Energy audits; LED systems; thermal insulation; smart monitoring; wider deployment of renewable solutions
& Alignment reinforcement gap
& Moderate \\

& Energy control and monitoring
& Weakly structured or insufficiently visible control systems
& Public performance indicators; strategic governance; systematic monitoring
& Alignment reinforcement gap
& Moderate \\

\addlinespace

\multirow{3}{*}{Human capital development}
& Training and employability
& CSR integration in curricula; future-oriented skills development
& Stronger alignment with green careers and sustainability competencies
& Alignment reinforcement gap
& Moderate \\

& Corporate partnerships
& Partnerships not systematically sustainability-oriented
& Responsible partnerships with certified companies and NGOs
& Support gap
& Moderate \\

& Sustainable student initiatives
& Existing sustainable entrepreneurship projects
& Dedicated incubators; sustainability-focused internships
& Support gap
& Moderate \\

\addlinespace

Responsible financial governance &
Budget transparency
& No clearly identified transparency mechanisms
& Dedicated sustainability budget; institutionalized green funds; public reporting; quantified impact indicators
& Institutional gap
& High \\

\addlinespace

\multirow{2}{*}{Responsible investment}
& Green budgeting
& No explicit green budget identified
& Formalized green funds; dedicated sustainability budget lines
& Strategic gap
& High \\

& Ecological investment policy
& Digitalization and resource optimization initiatives
& Strengthened ecological procurement; diversified sustainable investments; impact monitoring
& Strategic gap
& Moderate \\

\end{longtable}

\subsection{Social Expectation--Perception Gap}

The analysis of the social dimension reveals a partial alignment between perceived initiatives and stakeholder expectations, yet with significant qualitative gaps. While respondents acknowledge the existence of inclusion efforts, student clubs, integration events, and community engagement activities, expectations tend to emphasize the need for greater institutionalization and structural reinforcement, particularly regarding psychological support services and anti-discrimination mechanisms. Perception gaps are mainly observed in the following themes: inclusion and equity, well-being and psychological support, and community engagement. The theme clubs and leisure activities is considered largely aligned with stakeholder expectations.

\subsubsection{Formalization Gap: Inclusion and Equity}

The comparison of perceived initiatives and stakeholder expectations regarding inclusion and equity reveals a primarily qualitative gap. Inclusion and equity constitute a central pillar of the social dimension. They are mainly reflected in integration events, awareness campaigns, and support mechanisms for students facing difficulties. This theme is the most visible and widely recognized by respondents. With 32 mentions, it represents the dominant component of perceived social sustainability. Among these actions, integration events are the most frequently cited (27 mentions), indicating that the institution prioritizes concrete and visible measures to promote diversity, inclusion, and cohesion on campus. These events serve as the main channel through which equity efforts are perceived.

In contrast, awareness-raising initiatives and support for students in difficulty are less frequently mentioned, suggesting that, while such mechanisms exist, they are less developed or less visible to the university community. Students acknowledge efforts promoting integration and diversity through events, student clubs, campaigns, and certain anti-discrimination initiatives.

However, stakeholder expectations reflect a desire for greater formalization and institutionalization. Respondents call for structured anti-harassment programs, permanent mechanisms to combat discrimination, regular diversity training, and increased student participation in social decision-making. The observed gap, therefore, does not stem from a lack of initiatives but rather from the need to consolidate and strengthen existing actions.

\subsubsection{Structured Gap: Student Well-being}

Student well-being encompasses initiatives aimed at improving quality of life on campus, including psychological support, wellness activities, and dedicated events. While these measures exist, they are less prominently perceived than those related to inclusion. Psychological support, mentioned eight times in perceptions, indicates that mental health is recognized as an important component of university well-being. However, the relatively low frequency of mentions suggests that these services could be better structured, reinforced, and made more accessible.

Stakeholder expectations, in contrast, highlight this sub-theme (mentioned 26 times), emphasizing the need for a formal counseling unit, a permanent support center, and easier access to qualified professionals, particularly during exam periods.

\begin{quote}
Establish a psychological counseling unit as a permanent support service, provide more training sessions and workshops for students, and implement more regular monitoring of academic stress. (Respondent 15, Student)
\end{quote}

Academic stress management also emerges as a priority expectation, with students requesting regular workshops, preventive strategies during high-pressure periods, and structured rather than ad hoc support. Broader aspects of campus life quality (including rest areas, clear communication of available services, and a balance between academic demands and personal well-being) reflect a desire for a more supportive environment that fosters overall development.

Overall, this gap can be characterized as a student well-being structuring gap, where initiatives exist but require stronger organization and accessibility.

\subsubsection{Formalization Gap: Community Engagement}

Community engagement, mentioned 16 times, is a relatively well-identified dimension, reflecting a positive perception of the institution’s involvement in socially and civically oriented actions. Analysis of this theme, through the sub-themes of volunteering and solidarity projects, reveals a gap in formalization between student perceptions and expectations. Respondents recognize a dynamic campus life, marked by active clubs, volunteering activities, fundraising efforts, and occasional solidarity events. However, these initiatives are largely perceived as episodic and dependent on the commitment of student organizations, rather than being part of a structured institutional strategy.

Stakeholder expectations emphasize the need for more structured and regular volunteering opportunities, based on sustainable partnerships with NGOs, recurring solidarity projects with measurable impact, and better coordination between the administration and student associations. Additionally, students seek greater involvement in decision-making regarding the institution’s social and community orientations, particularly in the design and governance of socially impactful projects.

Thus, the identified gap does not indicate an absence of initiatives but rather underscores the need for formalization of community engagement at the institutional level.

\subsubsection{Alignment Reinforcement Gap: Culture \& Leisure}

Culture and leisure encompass cultural activities, sports events, and student clubs, all of which actively contribute to campus life and social cohesion. While these initiatives are recognized, they are cited less frequently compared to other social sub-themes, suggesting potential for further development in promoting student participation and associative vitality.

Unlike other social sub-themes such as psychological support, anti-discrimination measures, and participation in decision-making, student clubs and leisure activities appear mainly in perceived initiatives and less in priority expectations. This indicates a high level of alignment between what is offered and what is expected. Students acknowledge a dynamic campus life, with cultural, sports, and integration activities, without expressing major demands for improvement. The gap in this area is therefore low or almost nonexistent, indicating that culture and leisure represent a stable and well-aligned element within the social dimension of university sustainability.

Overall, social sustainability is perceived positively, with strong recognition of actions related to inclusion. However, the analysis highlights an uneven structuring of initiatives, with some dimensions relying more on ad hoc and visible actions rather than fully consolidated institutional mechanisms.

\renewcommand{\arraystretch}{1.2}
\footnotesize

\renewcommand{\arraystretch}{1.2}
\footnotesize

\begin{longtable}{
>{\raggedright\arraybackslash}p{2.4cm}
>{\raggedright\arraybackslash}p{2.8cm}
>{\raggedright\arraybackslash}p{3.5cm}
>{\raggedright\arraybackslash}p{3.6cm}
>{\raggedright\arraybackslash}p{3.0cm}
>{\centering\arraybackslash}p{1.2cm}
}

\caption{Social Expectation--Perception Gap Analysis}
\label{tab:social_gap} \\

\toprule
\textbf{Dimension} &
\textbf{Sub-theme} &
\textbf{Perception (existing initiatives)} &
\textbf{Priority expectations} &
\textbf{Nature of the gap} &
\textbf{Level} \\
\midrule
\endfirsthead

\toprule
\textbf{Dimension} &
\textbf{Sub-theme} &
\textbf{Perception (existing initiatives)} &
\textbf{Priority expectations} &
\textbf{Nature of the gap} &
\textbf{Level} \\
\midrule
\endhead

\midrule
\multicolumn{6}{r}{\textit{Continued on next page}} \\
\midrule
\endfoot

\bottomrule
\endlastfoot

%================================================
\multicolumn{6}{l}{\textbf{Inclusion \& Equity}} \\
\midrule

Inclusion \& Equity
& Anti-discrimination and anti-harassment programs
& Inclusion initiatives, awareness events, inclusive student clubs, equity promotion, occasional campaigns
& Stronger anti-discrimination programs, formal anti-harassment policies
& Formalization gap: initiatives exist but lack institutional reinforcement
& Moderate \\

& Support for vulnerable students
& Support for students in academic difficulty
& Increased accessibility and personalized academic support
& Gap in personalized academic support and accessibility
& Moderate \\

\addlinespace[0.8em]

%================================================
\multicolumn{6}{l}{\textbf{Student Well-being}} \\
\midrule

Student well-being
& Listening and psychological support
& Some psychological services, well-being activities, support for students in difficulty
& Permanent counseling unit, listening center, more psychologists, structured stress-management programs
& Structured gap: demand for institutionalized and continuous services
& High \\

& Well-being and stress management
& Well-being events and student activities
& Regular stress-management workshops, rest spaces, academic stress monitoring
& Regularity gap
& Moderate \\

& Student life and quality of life
& Occasional and informal activities
& Enhanced follow-up during exams, regular mental health events
& Structured gap
& Moderate \\

\addlinespace[0.8em]

%================================================
\multicolumn{6}{l}{\textbf{Community Involvement}} \\
\midrule

Community involvement
& Solidarity projects
& Integration events, solidarity projects
& Expanded social projects, NGO partnerships, broader engagement opportunities
& Alignment reinforcement gap: convergence exists but expectations call for scaling and diversification
& Moderate \\

& Volunteering
& Volunteer initiatives
& Structured volunteer programs with recognition and broader student participation
& Formalization gap
& Moderate \\

& Student involvement in decision-making
& Limited student representation in committees
& Greater involvement in institutional decision-making processes
& Formalization gap
& Moderate \\

\addlinespace[0.8em]

%================================================
\multicolumn{6}{l}{\textbf{Clubs \& Leisure}} \\
\midrule

Clubs \& Leisure
& Cultural and sports clubs
& Cultural and sports clubs, student events, leisure activities
& Rest spaces, reinforcement and diversification of activities
& Alignment reinforcement gap
& Low \\

\end{longtable}

\subsection{Environmental Expectation--Perception Gap}

The analysis of the interviews highlights a strong thematic convergence around three main axes: 
(1) environmental management and reduction of ecological footprint, 
(2) green space planning, and 
(3) environmental awareness and culture. 
These themes appear both in the initiatives implemented and in stakeholder expectations. However, while the analysis remains focused on the same overarching themes, some gaps emerge at the sub-theme level.

\subsubsection{Alignment Reinforcement Gap: Environmental Management}

Environmental management represents the most visible set of initiatives within the environmental dimension and is the most frequently mentioned by the university community. Perceived initiatives primarily include waste sorting, recycling bins, reduced printing and paper use, and responsible water management. These actions reflect a concrete operational orientation, focusing on tangible daily practices that are easy to implement.

Despite their presence, these initiatives are often partial, inconsistently deployed, and sporadic rather than integrated into a structured institutional strategy. For example, waste management is frequently limited to a few recycling bins and reduced printing. Students express more ambitious expectations, such as the installation of additional recycling bins, the development of a structured and institutionalized recycling system, the introduction of composting programs, and improved organization of waste collection. Notably, the absence of a ``zero-plastic'' policy highlights a major gap: although plastic reduction is recognized as a strategic environmental priority, it has not yet been fully incorporated into current practices.

Water management and optimization also emerge as strong expectations. Stakeholders propose concrete measures such as water-saving devices, rainwater harvesting systems, and the provision of water fountains. While some actions already exist and are spontaneously acknowledged, they remain largely invisible and insufficiently communicated, limiting their perceived impact.

Overall, although these initiatives are in place and recognized by the university community, they are perceived as limited in both intensity and reach. The main gap, therefore, does not reflect a lack of action but rather the need to strengthen the scale and consistency of environmental management initiatives.

\subsubsection{Structured Gap: Sustainable Campus Planning}

Sustainable campus planning constitutes the second identified theme. Respondents primarily highlight the presence and maintenance of green spaces, which are perceived as a positive element contributing to an improved campus environment. Sustainable mobility, notably through references to public transport and carpooling, is also mentioned, although these initiatives appear limited and poorly formalized.

Analysis reveals a partial alignment between perceptions and expectations. While green spaces are recognized and initiatives in this area are generally aligned with stakeholder perceptions, sustainable mobility emerges as a higher-priority expectation. Stakeholders call for more structured and organized actions to support public transport, carpooling, and other sustainable mobility solutions.

This shift in emphasis indicates that, although the landscaping dimension is acknowledged, the main structured gap lies in mobility management. The community expects a more integrated and functional vision of a sustainable campus, where mobility, spatial organization, and overall coherence are central. Current initiatives, while visible, remain fragmented and insufficiently structured, particularly regarding mobility programs.

Overall, sustainable campus planning is well established in its green space component, but stakeholders’ expectations highlight the need to expand mobility initiatives, making it a key organizational priority.

\subsubsection{Formalization Gap: Environmental Awareness and Culture}

The analysis of the theme environmental awareness and culture reveals a moderate gap between perceived initiatives and stakeholder expectations. Regarding perceived initiatives, respondents mention the existence of awareness campaigns and, to a lesser extent, student involvement in green projects. However, the relatively low number of occurrences suggests that these initiatives are neither systematic nor regular. Environmental culture appears to be still under development, relying primarily on ad hoc actions rather than on a continuous, integrated institutional dynamic. Awareness efforts are largely associated with occasional activities and student participation in green projects, making them present but limited in scope and structure.

In contrast, stakeholder expectations place strong emphasis on awareness campaigns, identifying them as a priority lever for promoting sustainability. This reflects a clear desire for the formalization of environmental communication initiatives beyond isolated or sporadic actions.

\begin{quote}
I would like to see more widespread waste sorting and improved waste management. I also hope that soft mobility options are more actively encouraged, and that students and staff are regularly made aware of sustainable development issues. Finally, I would like to see more collaborative sustainability projects involving both students and teachers. (Respondent 13, Student)
\end{quote}

The observed gap can therefore be characterized as a formalization gap. Perceived initiatives largely represent individual or isolated commitments, whereas stakeholder expectations reflect a comprehensive and integrated vision, in which environmental awareness is embedded within a coherent, visible, and continuous institutional strategy.

Overall, while environmental culture is recognized as being initiated, the university community expects it to evolve into a truly structuring pillar of a sustainable campus, linking communication, participation, and infrastructure transformation.

\begin{longtable}{p{2.8cm} p{2.5cm} p{3.3cm} p{3.3cm} p{2.5cm} p{1.2cm}}

\caption{Environmental Expectation--Perception Gap Analysis} \\

\toprule
\textbf{Dimension} & \textbf{Sub-theme} &
\textbf{Perception (existing initiatives)} &
\textbf{Priority expectations} &
\textbf{Nature of the Gap} &
\textbf{Gap level} \\
\midrule
\endfirsthead

\toprule
\textbf{Dimension} & \textbf{Sub-theme} &
\textbf{Perception (existing initiatives)} &
\textbf{Priority expectations} &
\textbf{Nature of the Gap} &
\textbf{Gap level} \\
\midrule
\endhead

% ==========================================================
% Environmental Management
% ==========================================================

\multicolumn{6}{l}{\textbf{Environmental Management}} \\
\midrule

Environmental Management
& Waste Management
& Selective sorting, few bins, paper recycling
& Structured system, more collection points, composting, effective management
& Alignment Reinforcement Gap
& Moderate \\

Environmental Management
& Plastic Reduction
& Rarely mentioned
& Zero-plastic policy, water fountains, elimination of plastic bottles
& Alignment Reinforcement Gap
& Moderate \\

Environmental Management
& Water Management
& Limited visibility
& Water-saving devices, rainwater harvesting systems
& Alignment Reinforcement Gap
& Moderate \\

Environmental Management
& Digitalization / Paper Reduction
& Well perceived
& Continue and reinforce
& Alignment Reinforcement Gap
& Low \\

\midrule
\addlinespace[0.6em]

% ==========================================================
% Sustainable Campus Planning
% ==========================================================

\multicolumn{6}{l}{\textbf{Sustainable Campus Planning}} \\
\midrule

Sustainable Campus Planning
& Green Space
& Maintained areas, plantings
& Further development, biodiversity enhancement
& Alignment Reinforcement Gap
& Low \\

Sustainable Campus Planning
& Sustainable Mobility
& Marginally mentioned
& Structured programs for public transport, carpooling, and other mobility solutions
& Structured Gap
& High \\

\midrule
\addlinespace[0.6em]

% ==========================================================
% Environmental Awareness and Culture
% ==========================================================

\multicolumn{6}{l}{\textbf{Environmental Awareness and Culture}} \\
\midrule

Environmental Awareness and Culture
& Awareness Campaigns
& Occasional campaigns
& Regular and structured campaigns integrated into the curriculum
& Alignment Reinforcement Gap
& Moderate \\

Environmental Awareness and Culture
& Student Involvement in Green Projects
& Limited participation
& Strengthened and structured engagement
& Alignment Reinforcement Gap
& Moderate \\

\bottomrule

\end{longtable}

\section{University Sustainability Index Formulation}

In order to formalize the proposed sustainability framework, university sustainability is modeled using latent variables representing stakeholder perception, stakeholder expectation, sustainability alignment, and overall institutional sustainability performance. Such constructs are not directly observable and must therefore be inferred from observable indicators through latent variable modeling techniques \citep{Bollen1989,Kaplan2009}. Latent variable modeling is particularly appropriate when dealing with multidimensional and abstract constructs such as sustainability \citep{Hair2019}.

\subsection{Definition of Latent Variables and Constructs}

The model distinguishes between exogenous latent variables, representing perceived and expected sustainability, and endogenous latent variables, representing sustainability gaps and the overall University Sustainability Index.

The vector of exogenous latent variables is defined as:

\begin{equation}
\boldsymbol{\xi} =
\begin{bmatrix}
P_e \\
P_s \\
P_{env} \\
E_e \\
E_s \\
E_{env}
\end{bmatrix}
\label{eq:xi}
\end{equation}

where:

\begin{itemize}
\item $P_e$: Perceived Economic Sustainability;
\item $P_s$: Perceived Social Sustainability;
\item $P_{env}$: Perceived Environmental Sustainability;
\item $E_e$: Expected Economic Sustainability;
\item $E_s$: Expected Social Sustainability;
\item $E_{env}$: Expected Environmental Sustainability.
\end{itemize}

These exogenous latent variables capture stakeholder-based evaluations from both perceptual and normative perspectives, consistent with expectation–perception models \citep{Parasuraman1988}.

The vector of endogenous latent variables is defined as:

\begin{equation}
\boldsymbol{\eta} =
\begin{bmatrix}
G_e \\
G_s \\
G_{env} \\
USI
\end{bmatrix}
\label{eq:eta}
\end{equation}

where:

\begin{itemize}
\item $G_e$: Economic Sustainability Gap;
\item $G_s$: Social Sustainability Gap;
\item $G_{env}$: Environmental Sustainability Gap.
\end{itemize}

The sustainability gap for each dimension is defined as:

\begin{align}
G_e &= E_e - P_e \label{eq:ge} \\
G_s &= E_s - P_s \label{eq:gs} \\
G_{env} &= E_{env} - P_{env} \label{eq:genv}
\end{align}

The University Sustainability Index (USI) is defined as a higher-order latent construct integrating the three sustainability gaps \citep{Bollen1989,Ringle2015}.

\begin{equation}
USI = \beta_1 G_e + \beta_2 G_s + \beta_3 G_{env} + \zeta
\label{eq:usi}
\end{equation}

where:

\begin{itemize}
\item $\beta_1, \beta_2, \beta_3$ are the contribution weights;
\item $\zeta$ is the structural disturbance term.
\end{itemize}

Composite index construction through weighted aggregation is widely used in sustainability assessment frameworks \citep{OECD2008,Nardo2005}.

\subsection{Interpretation of the Index}

The proposed USI reflects institutional sustainability as an alignment-based construct rather than a purely operational performance measure. Alignment-based models emphasize congruence between organizational actions and stakeholder expectations \citep{Suchman1995}.

Higher values indicate:

\begin{itemize}
\item Strong alignment between expectations and perceptions;
\item Effective institutional sustainability strategies;
\item Higher perceived sustainability legitimacy.
\end{itemize}

\subsection{Alignment-Based Reformulation}

Since sustainability performance increases when the gap decreases:

\begin{equation}
USI = 1 - \left( \beta_e |G_e| + \beta_s |G_s| + \beta_{env} |G_{env}| \right)
\label{eq:alignment}
\end{equation}

\subsection{Normalized Index Formulation}

For cross-institutional comparability:

\begin{equation}
USI^* = \frac{USI - USI_{\min}}{USI_{\max} - USI_{\min}}
\label{eq:normalized}
\end{equation}

Index normalization techniques are commonly employed in composite indicator construction \citep{OECD2008,Nardo2005}.

\section{Discussion}

This study demonstrates that the sustainable performance of universities cannot be fully captured by technical or operational indicators alone. Qualitative analyses reveal a \textit{Sustainability Expectation--Perception Gap}, defined as a systematic divergence between stakeholders’ perceptions and the initiatives implemented by institutions. This gap highlights limitations in visibility, communication, and strategic alignment that traditional performance measures fail to capture. Recognizing this gap emphasizes the importance of integrating stakeholders’ perspectives into sustainability assessments. Building on these insights, the study proposes a future integrative model incorporating perception-based variables into sustainable performance indices, enabling a more holistic and behaviorally informed evaluation across economic, social, and environmental dimensions.

\subsection{Economic Dimension}

The economic dimension of sustainability at ESB is characterized by effective operational implementation, primarily through digitalization and resource optimization. Digitalization dominates stakeholder responses and reflects a pragmatic understanding of economic sustainability, focused on immediately visible resource savings, particularly paper reduction. 

In contrast, energy optimization initiatives, although present, are less frequently mentioned and less perceptible. Concrete actions such as intelligent room management or reducing energy waste are identified, while more technical or strategic measures (e.g., solar panels, energy control systems) remain largely unrecognized. This situation appears to reflect not a lack of initiatives, but rather a deficit in visibility and communication, particularly regarding investments in sustainable infrastructure and their long-term impact.

Human capital development emerges as a more abstract dimension, primarily driven by faculty and less integrated into students’ perception of sustainability. The observed gap does not stem mainly from a lack of actions but from insufficient strategic valorization: weak communication, limited budgetary transparency, absence of visible quantitative indicators, and the lack of a formalized and structured framework (e.g., green budgeting, CSR reporting).

\subsection{Social Dimension}

The analysis of the social dimension reveals a mix of recognized initiatives and significant gaps. Although respondents acknowledge efforts in inclusion, student clubs, integration events, and community engagement activities, expectations emphasize the need for greater institutionalization and structural reinforcement, particularly regarding psychological support services and anti-discrimination mechanisms.

The most pronounced gap concerns mental health support. While some initiatives are perceived, stakeholders call for permanent, structured, and more accessible psychological services. Expectations also indicate a shift from occasional event-based participation toward greater involvement in institutional decision-making, reflecting a demand for participatory governance.

Overall, the findings suggest that the university’s social sustainability is perceived as active but insufficiently formalized, highlighting a transition from symbolic engagement toward structured and institutionalized social responsibility.

\subsection{Environmental Dimension}

The environmental dimension shows partial alignment between perceived initiatives and stakeholder expectations. Waste, plastic, and water management, as well as environmental awareness and student engagement, are moderately implemented but require stronger structuring and visibility. 

Sustainable mobility presents the most significant gap, reflecting a clear demand for formalized and structured mobility programs. Overall, the university’s environmental sustainability appears active yet insufficiently institutionalized, highlighting the need for reinforcement, strategic integration, and consistent stakeholder involvement.

\section{Limitations}

This study has several limitations. First, it relies primarily on qualitative interviews and the perceptions of internal stakeholders, which are inherently context-specific. The perspectives of external stakeholders were not included, although they could provide valuable additional insights into institutional sustainability performance.

Second, the findings may not be directly transferable to other universities or contexts; they should be interpreted in terms of contextual applicability rather than broad generalization. Finally, the study focuses on identifying and conceptualizing the Sustainability Expectation--Perception Gap without empirically testing or operationalizing it within a quantitative framework. Future research could address these limitations by incorporating external stakeholder perspectives and examining the applicability of the model across diverse institutional contexts.

\section{Theoretical Implications}

These findings contribute to the literature by demonstrating that sustainable performance in universities is not solely determined by technical measures but also by stakeholder perception and alignment. The Sustainability Expectation--Perception Gap underscores the importance of integrating behavioral and perceptual dimensions into sustainability assessment models, thereby extending traditional operational approaches. 

By conceptualizing sustainability as an alignment-based construct, this study advances stakeholder theory within the higher education context and offers a more comprehensive analytical framework for evaluating institutional sustainability performance.

\section{Managerial Implications}

From a managerial perspective, universities can use these insights to recognize and reduce the Sustainability Expectation--Perception Gap. Doing so may enhance the effectiveness of sustainability initiatives, improve stakeholder satisfaction, and foster more active engagement among students and staff.

This alignment-based approach provides decision-makers with a strategic framework for reconciling institutional initiatives with stakeholder expectations. By moving beyond purely technical indicators, universities can strengthen both the legitimacy and long-term sustainability of their actions.

%------------------------------------------------
% BIBLIOGRAPHY
%------------------------------------------------
\bibliographystyle{apacite}
\bibliography{references}

\end{document}